\section{Potentials}

Throughout, $(E, \mathcal{E})$ is a measurable space, $S$ a set of sites,
$(\Cfg, \mathcal{F}) = (E, \mathcal{E})^S$, $\Fin$ the set of finite subsets of $S$, and
$\beta \in \R$ an inverse temperature, so that the Boltzmann factor is
$e^{-\beta H^\Phi_\Lambda}$ and Georgii's statements are the case $\beta = 1$. Potentials are
indexed by all of $\Finset(S)$; the normalization $\Phi_\emptyset = 0$ is part of membership in
$\Bspace$.

\begin{definition}[Interaction potential, Georgii (2.2)(i)]
    \label{def:pot-potential}
    \uses{def:cylinder-event}
    \lean{Potential, Potential.IsPotential}
    \leanok

    A \textbf{potential} is a family $\Phi = (\Phi_A)_{A \in \Fin}$ of functions
    $\Phi_A : \Cfg \to \R$. It is an \textbf{interaction potential} if for each $A \in \Fin$
    the function $\Phi_A$ is $\mathcal{F}_A$-measurable.
\end{definition}

\begin{lemma}[An interaction term is local]
    \label{lem:pot-depends-on}
    \uses{def:pot-potential}
    \lean{Potential.IsPotential.dependsOn, Potential.IsPotential.eq_of_eqOn}
    \leanok

    If $\Phi$ is an interaction potential and $A \in \Fin$, then $\Phi_A(\eta) = \Phi_A(\zeta)$
    whenever $\eta_i = \zeta_i$ for all $i \in A$; equivalently, $\Phi_A$ is a function on
    $(E^A, \mathcal{E}^A)$.
\end{lemma}
\begin{proof}
    \leanok

    A real-valued $\mathcal{F}_A$-measurable function is constant on the fibres of
    $\eta \mapsto \eta_A$, since no set of $\mathcal{F}_A$ separates two configurations
    agreeing on $A$.
\end{proof}

\begin{definition}[Hamiltonian, Georgii (2.1), (2.2)(ii) and (2.3)]
    \label{def:pot-hamiltonian}
    \uses{def:pot-potential}
    \lean{Potential.IsSummable, Potential.hamiltonianTerms, Potential.hamiltonian}
    \leanok

    A potential $\Phi$ is \textbf{summable} if for all $\Lambda \in \Fin$ and $\omega \in \Cfg$
    the series
    \[ H_\Lambda^\Phi(\omega) = \sum_{A \in \Fin,\ A \inter \Lambda \neq \emptyset}
        \Phi_A(\omega) \]
    exists in the sense of Convention (2.1): the net
    \[ \Bigl( \sum_{A \inter \Lambda \neq \emptyset,\ A \subset \Delta} \Phi_A(\omega)
      \Bigr)_{\Delta \in \Fin} , \]
    indexed by the finite volumes $\Delta$ ordered by inclusion, converges. Its limit
    $H^\Phi_\Lambda$ is the \textbf{Hamiltonian} in $\Lambda$; its summands are the $\Phi_A$,
    with $\Phi_A$ replaced by $0$ when $A \inter \Lambda = \emptyset$.
\end{definition}

\begin{definition}[Boltzmann factor, Georgii (2.4)]
    \label{def:pot-boltzmann-factor}
    \uses{def:pot-hamiltonian}
    \lean{Potential.boltzmannFactor}
    \leanok

    $h_\Lambda^\Phi = \exp[-\beta H_\Lambda^\Phi]$, regarded as a $[0, \infty]$-valued density.
    It is everywhere strictly positive and finite.
\end{definition}

\begin{lemma}[Georgii (2.6)]
    \label{lem:pot-hamiltonian-diff}
    \uses{def:pot-hamiltonian, lem:pot-depends-on}
    \lean{Potential.dependsOn_hamiltonian_sub, Potential.hamiltonian_sub_eq_of_subset_eqOn_compl}
    \leanok

    Let $\Phi$ be a summable interaction potential and $\emptyset \neq \Lambda \subset \Delta
    \in \Fin$. Then
    \[ H^\Phi_\Delta - H^\Phi_\Lambda = \sum_{A \inter \Delta \neq \emptyset,\
        A \inter \Lambda = \emptyset} \Phi_A , \]
    and this difference depends only on the coordinates in $S \setminus \Lambda$:
    $H^\Phi_\Lambda(\eta) - H^\Phi_\Lambda(\zeta) = H^\Phi_\Delta(\eta) -
    H^\Phi_\Delta(\zeta)$ whenever $\eta_{S \setminus \Lambda} = \zeta_{S \setminus \Lambda}$.
    For countable $S$ it is therefore $\mathcal{F}_{S \setminus \Lambda}$-measurable.
\end{lemma}
\begin{proof}
    \leanok

    Each term of the difference series is indexed by an $A$ disjoint from $\Lambda$, and such a
    $\Phi_A$ depends only on the coordinates outside $\Lambda$ by
    Lemma~\ref{lem:pot-depends-on}. The property passes to finite partial sums and to the limit
    along the summation net.
\end{proof}

\begin{proposition}[Georgii, Proposition (2.5)]
    \label{prop:pot-premodification}
    \uses{def:premodif, def:pot-boltzmann-factor, lem:pot-hamiltonian-diff}
    \lean{Potential.isPremodifier_boltzmannFactor, Potential.boltzmannFactor_pos}
    \leanok

    Let $S$ be countable and $\Phi$ a summable interaction potential. Then
    $(h_\Lambda^\Phi)_{\Lambda \in \Fin}$ is a pre-modification: each $h^\Phi_\Lambda$ is
    measurable, and $h_\Delta^\Phi(\zeta) h_\Lambda^\Phi(\eta) = h_\Lambda^\Phi(\zeta)
    h_\Delta^\Phi(\eta)$ whenever $\Lambda \subset \Delta$ and $\zeta_{S \setminus \Lambda} =
    \eta_{S \setminus \Lambda}$. It is everywhere strictly positive.
\end{proposition}
\begin{proof}
    \leanok

    Measurability: $H^\Phi_\Lambda$ is a pointwise limit of finite sums of measurable
    functions. For $\emptyset \neq \Lambda \subset \Delta$ and $\zeta_{S \setminus \Lambda} =
    \eta_{S \setminus \Lambda}$, Lemma~\ref{lem:pot-hamiltonian-diff} gives
    $H_\Lambda^\Phi(\zeta) + H_\Delta^\Phi(\eta) = H_\Delta^\Phi(\zeta) +
    H_\Lambda^\Phi(\eta)$; apply $x \mapsto e^{-\beta x}$. Positivity is $e^{x} > 0$.
\end{proof}

\begin{definition}[Partition function and $\lambda$-admissibility, Georgii (2.7)--(2.8)]
    \label{def:pot-partition-function}
    \uses{def:pot-boltzmann-factor, def:isssd}
    \lean{Specification.sigmaFiniteLambdaZ, Specification.IsSigmaFiniteLambdaAdmissible,
      Potential.sigmaFinitePartitionFunction, Potential.partitionFunction}
    \leanok

    Fix a $\sigma$-finite nonzero a priori measure $\lambda \in \Meas(E, \mathcal{E})$. A
    potential $\Phi$ is \textbf{$\lambda$-admissible} if the partition function
    \[ Z_\Lambda^\Phi(\omega) = \lambda_\Lambda h^\Phi_\Lambda(\omega)
       = \int \lambda^\Lambda(\mathrm{d}\zeta)\,
         \exp[-\beta H^\Phi_\Lambda(\juxtBy{\omega}(\zeta))] \]
    is finite for all $\Lambda \in \Fin$ and $\omega \in \Cfg$. Since $h^\Phi_\Lambda > 0$, an
    admissible $\Phi$ has $Z^\Phi_\Lambda > 0$ as well, so
    $\rho^\Phi_\Lambda = h^\Phi_\Lambda / Z^\Phi_\Lambda$ is defined.
\end{definition}

\begin{definition}[Gibbsian specification, Georgii Definition (2.9)]
    \label{def:pot-gibbs-spec}
    \uses{def:pot-partition-function, prop:pot-premodification, lem:premodif-modif, def:modif, def:spec, def:isssd-spec}
    \lean{Potential.gibbsSpecificationOfSigmaFiniteAdmissible,
      Potential.gibbsSpecificationOfAbsolutelySummable, Potential.gibbsSpecification}
    \leanok

    For a $\lambda$-admissible potential $\Phi$, the \textbf{Gibbs distribution} in $\Lambda$
    with boundary condition $\omega_{S \setminus \Lambda}$ is
    \[ \gamma^\Phi_\Lambda(A \mid \omega) = Z^\Phi_\Lambda(\omega)^{-1}
        \int \lambda^\Lambda(\mathrm{d}\zeta)\,
        \exp[-\beta H^\Phi_\Lambda(\juxtBy{\omega}(\zeta))]\,
        1_A(\juxtBy{\omega}(\zeta)) , \]
    and $\gamma^\Phi = \rho^\Phi \lambda$ is the \textbf{Gibbsian specification} for $\Phi$ and
    $\lambda$. It is a specification in each of the following cases: $S$ countable, $\Phi$ a
    summable interaction potential, $\lambda$ $\sigma$-finite nonzero and $\Phi$
    $\lambda$-admissible; $S$ countable, $\Phi$ an absolutely summable interaction potential and
    $\lambda$ a probability measure, admissibility holding by
    Corollary~\ref{cor:pot-admissible}; $S$ arbitrary, $\Phi$ a finite range interaction
    potential, $\lambda$ a probability measure and $Z^\Phi_\Lambda(\omega) < \infty$ for all
    $\Lambda, \omega$.
\end{definition}

\begin{remark}[Gibbs measures and phase transitions, Georgii (2.9)--(2.10)]
    \label{rmk:pot-gibbs-measures}
    \uses{def:gibbs-meas, def:pot-gibbs-spec}

    A random field $\mu \in \Gibbs{\Phi} := \Gibbs{\gamma^\Phi}$ is a \textbf{Gibbs
    measure} (a DLR state) for $\Phi$ and $\lambda$, and $\Phi$ \textbf{exhibits a phase
    transition} if $|\Gibbs{\Phi}| > 1$.
\end{remark}

\begin{definition}[Absolutely summable potentials, Georgii (2.11)--(2.12)]
    \label{def:pot-abs-summable}
    \uses{def:pot-potential}
    \lean{Potential.normAt, Potential.IsAbsolutelySummable}
    \leanok

    $\Phi$ is \textbf{absolutely summable} if for all $i \in S$
    \[ \|\Phi\|_i = \sum_{A \in \Fin,\ A \ni i} \|\Phi_A\| < \infty , \]
    where $\|\Phi_A\| = \sup_\eta |\Phi_A(\eta)|$ and the unordered sum is taken in
    $[0, \infty]$.
\end{definition}

\begin{lemma}[Georgii (2.11) $\Rightarrow$ (2.2)(ii)]
    \label{lem:pot-abs-summable-summable}
    \uses{def:pot-abs-summable, def:pot-hamiltonian}
    \lean{Potential.summable_hamiltonianTerms, Potential.IsAbsolutelySummable.isSummable,
      Potential.hamiltonian_eq_tsum}
    \leanok

    An absolutely summable potential is summable, and its Hamiltonian is the unconditional
    sum $H_\Lambda^\Phi(\eta) = \sum_A \Phi_A(\eta) 1_{A \inter \Lambda \neq \emptyset}$.
\end{lemma}
\begin{proof}
    \leanok

    The summand indexed by $A$ is bounded by $\|\Phi_A\| 1_{A \inter \Lambda \neq \emptyset}$,
    and $\sum_A \|\Phi_A\| 1_{A \inter \Lambda \neq \emptyset} \leq \sum_{i \in \Lambda}
    \|\Phi\|_i < \infty$, because an $A$ meeting $\Lambda$ contains some $i \in \Lambda$.
    Unconditional summability implies convergence of the net of Convention (2.1).
\end{proof}

\begin{remark}[Uniform convergence, Georgii (2.11) and (2.13)]
    \label{rmk:pot-uniform-convergence}
    \uses{lem:pot-abs-summable-summable}

    For an absolutely summable $\Phi$, the truncated Hamiltonians $H^\Phi_{\Lambda, \Delta} =
    \sum_{A \subset \Delta,\ A \inter \Lambda \neq \emptyset} \Phi_A$ of (2.13) converge to
    $H^\Phi_\Lambda$ uniformly in $\omega$ as $\Delta \uparrow S$.
\end{remark}

\begin{theorem}[Georgii (2.14)]
    \label{thm:pot-hamiltonian-bound}
    \uses{def:pot-abs-summable, def:pot-hamiltonian}
    \lean{Potential.enorm_hamiltonian_le, Potential.iSup_enorm_hamiltonian_le,
      Potential.abs_hamiltonian_le}
    \leanok

    For every absolutely summable $\Phi$ and $\Lambda \in \Fin$,
    \[ \sup_{\eta \in \Cfg} \bigl|H^\Phi_\Lambda(\eta)\bigr|
       \leq \sum_{i \in \Lambda} \|\Phi\|_i . \]
\end{theorem}
\begin{proof}
    \uses{lem:pot-abs-summable-summable}
    \leanok

    By Lemma~\ref{lem:pot-abs-summable-summable} the Hamiltonian is the unconditional sum of
    its terms, so the triangle inequality bounds it by
    $\sum_A \|\Phi_A\| 1_{A \inter \Lambda \neq \emptyset} \leq \sum_{i \in \Lambda}
    \|\Phi\|_i$.
\end{proof}

\begin{corollary}[Georgii (2.14) $\Rightarrow$ $\lambda$-admissibility]
    \label{cor:pot-admissible}
    \uses{thm:pot-hamiltonian-bound, def:pot-partition-function}
    \lean{Potential.isRelAdmissible_boltzmannFactor,
      Potential.isPremodifierAdmissible_boltzmannFactor,
      Potential.isSigmaFiniteLambdaAdmissible_boltzmannFactor}
    \leanok

    Let $\Phi$ be absolutely summable and $\beta \in \R$. For every specification $\gamma$,
    with no assumption on $S$, the partition functions $Z^\Phi_\Lambda = \gamma_\Lambda
    h^\Phi_\Lambda$ obey, for all $\Lambda, \omega$,
    \[ e^{-|\beta| \sum_{i \in \Lambda}\|\Phi\|_i} \ \leq\ Z^\Phi_\Lambda(\omega)\ \leq\
       e^{|\beta| \sum_{i \in \Lambda}\|\Phi\|_i} , \]
    so they are nonzero and finite, and $\rho^\Phi_\Lambda = h^\Phi_\Lambda / Z^\Phi_\Lambda
    \leq e^{2|\beta| \sum_{i \in \Lambda}\|\Phi\|_i}$. In particular $\Phi$ is
    $\lambda$-admissible for every finite nonzero $\lambda$.
\end{corollary}
\begin{proof}
    \leanok

    By Theorem~\ref{thm:pot-hamiltonian-bound}, $|H^\Phi_\Lambda| \leq \sum_{i \in \Lambda}
    \|\Phi\|_i$ uniformly in the configuration, so $h^\Phi_\Lambda$ lies between the constants
    $e^{\mp |\beta| \sum_{i \in \Lambda}\|\Phi\|_i}$. Integrating against a probability kernel
    preserves the constants, giving both bounds on $Z$; dividing gives the bound on
    $\rho^\Phi_\Lambda$. For finite nonzero $\lambda$, normalize $\lambda$ and take $\gamma$
    the $\lambda$-specification.
\end{proof}

\begin{definition}[The space $\Bspace$, Georgii (2.11)]
    \label{def:pot-bspace}
    \uses{def:pot-abs-summable}
    \lean{Potential.absolutelySummable, Potential.BSpace, Potential.seminormAt,
      Potential.seminormFamily}
    \leanok

    $\Bspace$ is the $\R$-linear space of all absolutely summable interaction potentials with
    $\Phi_\emptyset = 0$, equipped with the family of seminorms $\|\cdot\|_i$, $i \in S$.
\end{definition}

\begin{theorem}[$\Bspace$ is a Fr\'echet space, Georgii (2.11)]
    \label{thm:pot-frechet}
    \uses{def:pot-bspace}
    \lean{Potential.withSeminorms_seminormFamily, Potential.tendsto_iff_tendsto_seminormAt,
      Potential.eq_zero_of_forall_seminormAt_eq_zero,
      Potential.instMetrizableSpaceSubtypeMemSubmoduleRealAbsolutelySummableOfCountable,
      Potential.instCompleteSpaceSubtypeMemSubmoduleRealAbsolutelySummableOfCountable}
    \leanok

    For any $S$, the seminorms $\|\cdot\|_i$ define a locally convex $T_1$ topology on
    $\Bspace$, and $\Phi_x \to \Phi$ in $\Bspace$ iff $\|\Phi_x - \Phi\|_i \to 0$ for every
    $i \in S$. For countable $S$ this topology is metrizable and complete: $\Bspace$ is a
    Fr\'echet space.
\end{theorem}
\begin{proof}
    \leanok

    Convergence in the topology generated by a seminorm family is seminorm convergence. The
    seminorms separate points: if $\|\Phi\|_i = 0$ for every $i$ then $\sup_\eta|\Phi_A(\eta)|
    = 0$ for every nonempty $A$, so $\Phi = 0$; hence $T_1$. For countable $S$ the family is
    countable, so the separated space is metrizable. Completeness: a Cauchy sequence is Cauchy
    in each $\|\cdot\|_i$, hence each $\Phi_A(\eta)$ converges (a single term is dominated by
    $\|\cdot\|_i$ for any $i \in A$); the limit is absolutely summable by lower semicontinuity
    of $\|\cdot\|_i$ under pointwise convergence, and the convergence holds in every seminorm.
\end{proof}

\begin{lemma}[Convergence in $\Bspace$ makes Hamiltonians converge uniformly]
    \label{lem:pot-hamiltonian-uniform}
    \uses{thm:pot-hamiltonian-bound, def:pot-bspace}
    \lean{Potential.hamiltonian_sub, Potential.abs_hamiltonian_sub_le,
      Potential.hamiltonianBound_coe, Potential.tendsto_hamiltonianBound_sub}
    \leanok

    The Hamiltonian is linear in the potential on $\Bspace$, and for $\Phi, \Psi \in \Bspace$
    \[ \bigl|H^\Phi_\Lambda(\eta) - H^\Psi_\Lambda(\eta)\bigr|
       \leq \sum_{i \in \Lambda} \|\Phi - \Psi\|_i \qquad (\eta \in \Cfg) . \]
    Consequently $\Phi_x \to \Phi$ in $\Bspace$ implies
    $\sup_\eta |H^{\Phi_x}_\Lambda(\eta) - H^{\Phi}_\Lambda(\eta)| \to 0$ for every
    $\Lambda \in \Fin$.
\end{lemma}
\begin{proof}
    \leanok

    Linearity holds termwise and passes to the unconditional sums of
    Lemma~\ref{lem:pot-abs-summable-summable}, so $H^{\Phi - \Psi}_\Lambda = H^\Phi_\Lambda -
    H^\Psi_\Lambda$. Apply Theorem~\ref{thm:pot-hamiltonian-bound} to $\Phi - \Psi$; the bound
    $\sum_{i \in \Lambda} \|\Phi - \Psi\|_i$ is a finite sum of seminorms, which tends to $0$
    along a net converging in $\Bspace$ by Theorem~\ref{thm:pot-frechet}.
\end{proof}

\begin{proposition}[Georgii, Proposition (4.19); stated here because it is proved from (2.14)]
    \label{prop:pot-gibbs-uniform}
    \uses{lem:pot-hamiltonian-uniform, def:pot-gibbs-spec, cor:pot-admissible}
    \lean{Potential.dist_action_gibbsSpecification_le,
      Potential.tendsto_dist_action_gibbsSpecification,
      Potential.BSpace.tendsto_dist_action_gibbsSpecification}
    \leanok

    Let $S$ be countable, $\lambda$ a probability measure, and $\Phi, \Psi \in \Bspace$ with
    $|H^\Psi_\Lambda - H^\Phi_\Lambda| \leq D$ pointwise. Then for every bounded measurable
    $f$,
    \[ \bigl\| \gamma^\Psi_\Lambda f - \gamma^\Phi_\Lambda f \bigr\|
       \leq 2\|f\| \bigl( e^{|\beta| D} - 1 \bigr) , \]
    uniformly in the boundary condition. Hence, for a net $(\Phi_x)$ in $\Bspace$ with
    $\Phi_x \to \Phi \in \Bspace$ along any filter, $\gamma^{\Phi_x}_\Lambda f \to
    \gamma^{\Phi}_\Lambda f$ uniformly, for every $\Lambda \in \Fin$ and bounded measurable
    $f$.
\end{proposition}
\begin{proof}
    \leanok

    Write $\gamma^\Phi_\Lambda f = \lambda_\Lambda(h^\Phi_\Lambda f) /
    \lambda_\Lambda(h^\Phi_\Lambda)$. The estimate $|e^x - e^y| \leq
    e^{y}(e^{|x-y|} - 1)$ gives $|h^\Psi_\Lambda - h^\Phi_\Lambda| \leq h^\Phi_\Lambda
    (e^{|\beta| D} - 1)$ pointwise, hence $|Z^\Psi_\Lambda - Z^\Phi_\Lambda| \leq
    Z^\Phi_\Lambda (e^{|\beta| D} - 1)$ and $|\lambda_\Lambda(h^\Psi_\Lambda f) -
    \lambda_\Lambda(h^\Phi_\Lambda f)| \leq \|f\| Z^\Phi_\Lambda(e^{|\beta| D} - 1)$.
    Splitting $A/Z' - B/Z = A(Z - Z')/(Z'Z) + (A-B)/Z$ and using
    $|\lambda_\Lambda(h^\Psi_\Lambda f)| \leq \|f\| Z^\Psi_\Lambda$ bounds each half by
    $\|f\|(e^{|\beta| D} - 1)$. The limit statement follows from
    Lemma~\ref{lem:pot-hamiltonian-uniform} and $e^t - 1 \to 0$ as $t \to 0$.
\end{proof}

\begin{definition}[Finite range potentials, Georgii (2.15)]
    \label{def:pot-finite-range}
    \uses{def:pot-potential}
    \lean{Potential.IsFiniteRange, Potential.interactingSupport,
      Potential.interactingHamiltonian}
    \leanok

    $\Phi$ is of \textbf{finite range} if for each $i \in S$ there is $\Delta_i \in \Fin$ such
    that $\Phi_A = 0$ whenever $i \in A$ and $A \not\subset \Delta_i$. For such a $\Phi$ and
    $\Lambda \in \Fin$, the set of $A$ meeting $\Lambda$ with $\Phi_A \neq 0$ is finite, and
    the Hamiltonian in $\Lambda$ is the corresponding finite sum.
\end{definition}

\begin{lemma}[Finite range potentials are summable]
    \label{lem:pot-finite-range-summable}
    \uses{def:pot-finite-range, def:pot-hamiltonian}
    \lean{Potential.IsFiniteRange.isSummable, Potential.hasSum_interactingHamiltonian,
      Potential.hamiltonian_eq_interactingHamiltonian}
    \leanok

    A finite range potential satisfies (2.2)(ii), and its Hamiltonian coincides with the finite
    sum of Definition~\ref{def:pot-finite-range}.
\end{lemma}
\begin{proof}
    \leanok

    The summands vanish off a finite set of supports, so the net of partial sums is eventually
    constant.
\end{proof}

\begin{remark}[Pair potentials, Georgii (2.16)]
    \label{rmk:pot-pair}
    \uses{def:pot-potential}

    A \textbf{pair potential} is a potential with $\Phi_A = 0$ whenever $|A| > 2$; the standard
    examples are $\Phi_{\set{i,j}}(\omega) = J(i,j)\varphi(\omega_i, \omega_j)$,
    $\Phi_{\set{i}}(\omega) = J(i,i)\psi(\omega_i)$, with $J$ and $\varphi$ symmetric and
    $\varphi, \psi$ measurable.
\end{remark}

\begin{definition}[Nearest-neighbour potentials, Georgii (2.17)]
    \label{def:pot-nearest-neighbour}
    \uses{rmk:pot-pair, def:pot-finite-range}
    \lean{Potential.nearestNeighbourPair}
    \leanok

    Let $S$ be the vertex set of a simple graph $G$ with edge set $B$. A potential is a
    \textbf{nearest-neighbour} (or Markov) potential if $\Phi_A = 0$ except when $A$ is a
    complete subgraph of $(S, B)$. For a spin observable $\sigma : E \to \R$, a coupling
    $J \in \R$ and a field $h \in \R$, the \textbf{nearest-neighbour pair potential} of
    $(G, J, h, \sigma)$ is
    \[ \Phi_{\set{i}}(\eta) = -h\,\sigma(\eta_i), \qquad
       \Phi_{\set{i,j}}(\eta) = -J\,\sigma(\eta_i)\sigma(\eta_j)
       \ \text{ for } \set{i,j} \in B, \qquad \Phi_A = 0 \text{ otherwise},\]
    with the coefficients of (3.13). The Ising models are the case
    $E = \set{-1, +1}$, $\sigma = \mathrm{id}$.
\end{definition}

\begin{lemma}[The nearest-neighbour pair potential lies in $\Bspace$]
    \label{lem:pot-nearest-neighbour-summable}
    \uses{def:pot-nearest-neighbour, def:pot-abs-summable, def:pot-finite-range}
    \lean{Potential.isPotential_nearestNeighbourPair,
      Potential.isFiniteRange_nearestNeighbourPair,
      Potential.isAbsolutelySummable_nearestNeighbourPair,
      Potential.abs_nearestNeighbourPair_apply_le}
    \leanok

    If $\sigma$ is measurable then the nearest-neighbour pair potential of $(G, J, h, \sigma)$
    satisfies (2.2)(i). If moreover $G$ is locally finite it has finite range, and if in
    addition $|\sigma| \leq c$ pointwise it is absolutely summable, with
    $|\Phi_A(\eta)| \leq |h||c| + |J| c^2$ for all $A, \eta$.
\end{lemma}
\begin{proof}
    \leanok

    Measurability: on each of the three branches the term is a constant multiple of a finite
    sum or product of coordinate evaluations at sites of $A$. Finite range: a nonzero support
    containing $i$ is either $\set{i}$ or an edge at $i$, so it lies in
    $\set{i} \union N(i)$, which is finite by local finiteness. Absolute summability: only the
    supports inside $\set{i} \union N(i)$ contribute to $\|\Phi\|_i$, each at most
    $|h||c| + |J|c^2$.
\end{proof}

\begin{definition}[Self-potentials, Georgii (2.18)]
    \label{def:pot-self-potential}
    \uses{def:pot-potential, def:pot-hamiltonian}
    \lean{Potential.oneBody, Potential.manyBody, Potential.oneBody_add_manyBody,
      Potential.hamiltonian_oneBody, Potential.hamiltonian_manyBody}
    \leanok

    A \textbf{self-potential} is a potential with $\Phi_A = 0$ unless $|A| = 1$. Any $\Phi$
    splits by cardinality as $\Phi = \Phi^{1} + \Phi^{\geq 2}$, where $\Phi^{1}$ keeps the
    volumes with $|A| \leq 1$ (its \textbf{self-potential part}) and $\Phi^{\geq 2}$ those with
    $|A| \geq 2$. Then $\Phi^{1}$ has finite range with
    $H^{\Phi^{1}}_\Lambda(\eta) = \sum_{i \in \Lambda} \Phi_{\set{i}}(\eta)$, and for summable
    $\Phi$, $H^{\Phi^{\geq 2}}_\Lambda = H^\Phi_\Lambda - \sum_{i \in \Lambda}
    \Phi_{\set{i}}$.
\end{definition}

\begin{proposition}[Absorbing the self-potential into the a priori measure, Georgii (2.18)]
    \label{prop:pot-self-absorption}
    \uses{def:pot-self-potential, def:pot-gibbs-spec}
    \lean{Potential.lambdaSpecification_eq_gibbsSpecificationFamily}
    \leanok

    Let $S$ be countable, $\lambda$ $\sigma$-finite and nonzero, $\Phi$ a summable
    $\lambda$-admissible interaction potential, and $\eta_0 \in \Cfg$ such that the recentred
    many-body part $\bigl(\Phi^{\geq 2}_A - \Phi^{\geq 2}_A(\eta_0)\bigr)_{A \in \Fin}$ is
    absolutely summable. Let $c_i \in (0, \infty)$ be such that
    $\nu_i = c_i^{-1} e^{-\beta \Phi_{\set{i}}}\lambda$ is a probability measure for each
    $i \in S$. Then the $\lambda$-specification $\gamma^\Phi$ equals the Gibbsian specification
    of the recentred many-body part over the family $(\nu_i)_{i \in S}$.
\end{proposition}
\begin{proof}
    \leanok

    By Definition~\ref{def:pot-self-potential} the Boltzmann factor factorizes as
    $h^\Phi_\Lambda(\sigma) = e^{-\beta H^{\Phi^{\geq 2}}_\Lambda(\eta_0)}\cdot
    h^{\mathrm{many, centred}}_\Lambda(\sigma) \cdot \prod_{i \in \Lambda}
    e^{-\beta \Phi_{\set{i}}(\sigma)}$. The last product is the density turning
    $\lambda^\Lambda$ into $\bigotimes_{i \in \Lambda}(c_i \nu_i)$, and the volume-dependent
    constant in front cancels in the normalization. So the two normalized kernels agree on
    every volume and boundary condition.
\end{proof}
